\chapter*{Mở đầu}
\addcontentsline{toc}{chapter}{Mở đầu}
\markboth{MỞ ĐẦU}{}

\section*{Lý do chọn đề tài}

Các ứng dụng doanh nghiệp ngày càng sử dụng mô hình ngôn ngữ lớn (Large
Language Model -- LLM) và kiến trúc sinh văn bản có truy hồi tri thức
(Retrieval-Augmented Generation -- RAG). Bên cạnh lợi ích về khai thác tài liệu,
các hệ thống này mở rộng bề mặt tấn công thông qua prompt injection, chỉ dẫn độc
hại trong tài liệu truy hồi, jailbreak và rò rỉ dữ liệu. Nội dung do người dùng
hoặc nguồn tài liệu không tin cậy kiểm soát có thể bị mô hình diễn giải như chỉ
dẫn có thẩm quyền nếu hệ thống không phân tách và kiểm tra phù hợp.

Đề tài lựa chọn hướng phòng vệ nhiều lớp bằng một cổng bảo mật đặt trước bước
sinh câu trả lời. Hướng tiếp cận này phù hợp với phạm vi thực tập vì có thể hiện
thực bằng các rule giải thích được, chạy ngoại tuyến và đánh giá bằng dữ liệu
tổng hợp mà không sử dụng dữ liệu doanh nghiệp thật hoặc API trả phí.

\section*{Mục tiêu}

Mục tiêu của đề tài là nghiên cứu và triển khai một proof-of-concept cho LLM
Security Gateway/Guardrail Proxy, gồm các chức năng chính:

\begin{itemize}
    \item kiểm tra prompt trước khi chuyển sang bước xử lý tiếp theo;
    \item kiểm tra và làm sạch context chunk trước khi đưa vào provider;
    \item kiểm tra dữ liệu đầu ra trước khi trả về phía gọi;
    \item ghi log quyết định có cấu trúc và che giấu chuỗi nhạy cảm;
    \item đánh giá ngoại tuyến trên benchmark red-team tổng hợp, có thể tái lập.
\end{itemize}

\section*{Đối tượng và phạm vi nghiên cứu}

Đối tượng nghiên cứu là các cơ chế guardrail cho prompt đầu vào, ngữ cảnh RAG
và đầu ra, cùng quy trình đánh giá quyết định của guard. Hệ thống đã triển khai
FastAPI gateway, ba guard dựa trên rule/regex, mock provider cục bộ, audit log,
bộ nạp dữ liệu tổng hợp và evaluation runner.

Phạm vi không bao gồm LLM thực, API trả phí, embedding, vector database, truy
hồi tương đồng hoặc pipeline RAG hoàn chỉnh. Context trong demo được truyền trực
tiếp bởi phía gọi. Đề tài không xử lý dữ liệu thật và không tuyên bố hệ thống
sẵn sàng cho môi trường sản xuất.

\section*{Phương pháp thực hiện}

Nhóm tiến hành theo các bước: khảo sát tài liệu và khung rủi ro; xây dựng yêu
cầu, kiến trúc và threat model; tạo corpus tổng hợp; hiện thực gateway và các
guard; kiểm thử; chạy đánh giá có kiểm soát; phân tích false negative; hiệu chỉnh
rule hẹp; và so sánh chế độ guarded với baseline không guard. Mọi số liệu trong
báo cáo được lấy từ artifact có thể tái tạo trong kho mã nguồn.

\section*{Bố cục báo cáo}

Chương 1 trình bày nền tảng LLM/RAG, prompt injection, rò rỉ dữ liệu và
guardrail. Chương 2 mô tả yêu cầu, kiến trúc và threat model. Chương 3 trình bày
implementation, bộ dữ liệu và các module. Chương 4 mô tả thiết lập đánh giá,
kết quả, baseline và phân tích lỗi. Phần kết luận tổng hợp đóng góp, hạn chế và
hướng phát triển.
